\documentclass{article}

\title{A Sample Document}
\author{Test Author}
\date{2024-01-15}

\begin{document}

\maketitle

\begin{abstract}
This is the abstract of the document. It provides a brief summary
of the content discussed in the following sections.
\end{abstract}

\section{Introduction}
\label{sec:intro}

This is the introductory paragraph. It contains some \textbf{bold text}
and some \emph{emphasized text} for testing inline formatting removal.

% This is a comment and should be ignored by the parser.

\subsection{Background}

The background of the study is important. See Section~\ref{sec:intro}
for more details.

\subsubsection{Historical Context}

A brief look at the historical context of this work.

\paragraph{Early Work}

Some early work was done in this field by \cite{knuth1984}.

\section{Methods}

\subsection{Data Collection}

We collected data using the following methods:

\begin{itemize}
  \item First method of data collection
  \item Second method with \textbf{bold} details
  \item Third method of data collection
\end{itemize}

The steps of our analysis were:

\begin{enumerate}
  \item Preprocess the raw data
  \item Apply the transformation algorithm
  \item Validate the results
\end{enumerate}

\section{Results}

\begin{table}[h]
\centering
\caption{Experimental Results}
\label{tab:results}
\begin{tabular}{|l|c|r|}
\hline
Name & Score & Rank \\
\hline
Alice & 95 & 1 \\
Bob & 87 & 2 \\
Carol & 92 & 3 \\
\hline
\end{tabular}
\end{table}

The relationship is described by the following equation:

\begin{equation}
\label{eq:main}
E = mc^2
\end{equation}

We can also write inline-style display math:

$$
F = ma
$$

And bracket-style display math:

\[
\sum_{i=1}^{n} x_i = S
\]

\section{Discussion}

\begin{figure}[h]
\centering
\includegraphics{diagram.png}
\caption{A diagram showing the architecture}
\label{fig:architecture}
\end{figure}

As shown in Figure~\ref{fig:architecture}, the system has multiple layers.

\footnote{This is an important footnote about the discussion.}

\subsection{Code Example}

\begin{verbatim}
fn main() {
    println!("Hello, world!");
}
\end{verbatim}

\begin{lstlisting}[language=Python]
def hello():
    print("Hello, world!")
\end{lstlisting}

\section{Conclusion}

In conclusion, we have demonstrated the key findings of this study.
See Equation~\ref{eq:main} and Table~\ref{tab:results} for details.

\end{document}
